% !TeX root = ../main.tex
% Fonte letta integralmente: Lezione 1.pdf, pp. 7--13.
% Gli esempi finiti e il diagramma di Venn sono aggiunte esplicative.
\chapter{Insiemi e relazioni}
\label{cap:insiemi}
\fonti{Lezione 1}

\section{Insiemi, elementi e insieme ambiente}
% Fonte: Lezione 1, pp.~7--8

Il concetto di \emph{insieme} viene assunto come primitivo: intuitivamente indica una raccolta di oggetti, detti \emph{elementi} o \emph{punti}. Un insieme $A$ è assegnato quando, per ogni oggetto considerato, la frase «$x$ appartiene ad $A$» ha un valore di verità. Scriviamo
\[
  x\in A\quad\text{se $x$ appartiene ad $A$},\qquad
  x\notin A\quad\text{se non vi appartiene}.
\]
\textbf{Uguaglianza e rappresentazione.}
\begin{itemize}
  \item Due insiemi sono uguali se e solo se hanno gli stessi elementi.
  \item L'ordine dell'elenco e le eventuali ripetizioni non modificano l'insieme.
\end{itemize}

\textbf{Descrizione per elenco o per proprietà.} Si può assegnare un insieme elencandone gli elementi oppure indicando una proprietà che li caratterizza. Per esempio,
\[
  A=\{1,3,5,7,9\}
   =\{x\in\N:0\le x\le10\ \land\ x\text{ è dispari}\}.
\]
In generale, una descrizione per proprietà ha la forma $A=\{x\in U:P(x)\}$: $U$ è l'\emph{insieme ambiente} e si selezionano gli elementi di $U$ che soddisfano $P$.

\textbf{Insieme vuoto.} Si assume l'esistenza di un insieme privo di elementi, detto \emph{insieme vuoto} e indicato con $\varnothing$. Esso è diverso dal singoletto $\{\varnothing\}$: quest'ultimo contiene un elemento, che è a sua volta l'insieme vuoto.

\section{Sottoinsiemi e insieme delle parti}
% Fonte: Lezione 1, p.~9

\begin{definizione}[Inclusione]
Un insieme $B$ è un \emph{sottoinsieme} di $A$ se ogni elemento di $B$ è anche un elemento di $A$. Si scrive $B\subseteq A$, ossia
\[
  B\subseteq A\quad\Longleftrightarrow\quad
  \forall x,\ (x\in B\Rightarrow x\in A).
\]
\end{definizione}

\begin{itemize}
  \item \textbf{Sottoinsiemi impropri.} $\varnothing$ e $A$ sono sempre sottoinsiemi di $A$. L'inclusione $\varnothing\subseteq A$ vale perché il vuoto non ha elementi che possano violarla.
  \item \textbf{Inclusione stretta.} Se $B\subseteq A$ e $B\ne A$, scriviamo $B\subsetneq A$.
\end{itemize}

\textbf{Insieme delle parti.} L'\emph{insieme delle parti} di $A$, indicato con $\mathcal P(A)$, è l'insieme di tutti i suoi sottoinsiemi:
\[
  \mathcal P(A)=\{B:B\subseteq A\}.
\]
\begin{esempio}
Se $A=\{1,2\}$, allora $\mathcal P(A)=\{\varnothing,\{1\},\{2\},\{1,2\}\}$. È utile distinguere l'appartenenza dall'inclusione: $1\in A$, mentre $\{1\}\subseteq A$ e $\{1\}\in\mathcal P(A)$.
\end{esempio}

\section{Operazioni tra insiemi}
% Fonte: Lezione 1, pp.~10--11

Fissiamo due insiemi $A$ e $B$ contenuti nello stesso ambiente $U$. Le operazioni fondamentali traducono le corrispondenti operazioni logiche.

\begin{definizione}[Unione, intersezione e differenza]
Si definiscono
\begin{align*}
  A\cup B&=\{x\in U:x\in A\lor x\in B\},\\
  A\cap B&=\{x\in U:x\in A\land x\in B\},\\
  A\setminus B&=\{x\in U:x\in A\land x\notin B\}.
\end{align*}
L'unione corrisponde all'appartenenza ad almeno uno dei due insiemi; l'intersezione all'appartenenza a entrambi; la differenza all'appartenenza ad $A$ ma non a $B$.
\end{definizione}

\begin{schema}
\begin{tikzpicture}[x=1cm,y=1cm]
  \begin{scope}
    \fill[coloreschema!20] (-.45,0) circle[radius=.85];
    \fill[coloreschema!20] (.45,0) circle[radius=.85];
    \draw (-.45,0) circle[radius=.85];
    \draw (.45,0) circle[radius=.85];
    \node at (-.8,.35) {$A$}; \node at (.8,.35) {$B$};
    \node at (0,-1.2) {$A\cup B$};
  \end{scope}
  \begin{scope}[xshift=4cm]
    \begin{scope}
      \clip (-.45,0) circle[radius=.85];
      \fill[coloreschema!20] (.45,0) circle[radius=.85];
    \end{scope}
    \draw (-.45,0) circle[radius=.85];
    \draw (.45,0) circle[radius=.85];
    \node at (-.8,.35) {$A$}; \node at (.8,.35) {$B$};
    \node at (0,-1.2) {$A\cap B$};
  \end{scope}
  \begin{scope}[xshift=8cm]
    \fill[coloreschema!20] (-.45,0) circle[radius=.85];
    \fill[white] (.45,0) circle[radius=.85];
    \draw (-.45,0) circle[radius=.85];
    \draw (.45,0) circle[radius=.85];
    \node at (-.8,.35) {$A$}; \node at (.8,.35) {$B$};
    \node at (0,-1.2) {$A\setminus B$};
  \end{scope}
\end{tikzpicture}
\par\smallskip La regione colorata rappresenta l'operazione indicata sotto ciascun disegno.
\end{schema}

\begin{esempio}
Con $A=\{1,2,3\}$ e $B=\{3,4\}$ si ha $A\cup B=\{1,2,3,4\}$, $A\cap B=\{3\}$ e $A\setminus B=\{1,2\}$. Scambiando i due insiemi si ottiene invece $B\setminus A=\{4\}$: la differenza dipende dall'ordine.
\end{esempio}

\textbf{Complementare e insieme ambiente.} Se $B\subseteq A$, l'insieme $A\setminus B$ è detto \emph{complementare di $B$ in $A$}. Quando l'ambiente è $U$, si scrive semplicemente $B^c=U\setminus B$. L'ambiente è essenziale: il complemento dello stesso insieme cambia se cambia $U$.

\section{Le leggi di De Morgan}
% Fonte: Lezione 1, p.~11

\begin{proposizione}[Leggi di De Morgan]
Per $A,B\subseteq U$, con i complementi calcolati nello stesso ambiente,
\[
  (A\cup B)^c=A^c\cap B^c,\qquad
  (A\cap B)^c=A^c\cup B^c.
\]
\end{proposizione}
% Evita che la sola frase conclusiva della seconda identità passi alla pagina seguente.
\enlargethispage{\baselineskip}
\begin{proof}
\textbf{Prima identità.} Per ogni $x\in U$,
\[
  \begin{aligned}
    x\in(A\cup B)^c
    &\iff x\notin A\cup B\\
    &\iff x\notin A\ \land\ x\notin B\\
    &\iff x\in A^c\cap B^c.
  \end{aligned}
\]
Le due condizioni di appartenenza sono equivalenti per ogni elemento dell'ambiente, quindi i due insiemi sono uguali.

\textbf{Seconda identità.} Analogamente, per ogni $x\in U$,
\[
  \begin{aligned}
    x\in(A\cap B)^c
    &\iff x\notin A\cap B\\
    &\iff x\notin A\ \lor\ x\notin B\\
    &\iff x\in A^c\cup B^c.
  \end{aligned}
\]
Si ottiene anche la seconda uguaglianza.
\end{proof}

Le formule esprimono a livello di insiemi le equivalenze logiche $\neg(P\lor Q)\Leftrightarrow(\neg P\land\neg Q)$ e $\neg(P\land Q)\Leftrightarrow(\neg P\lor\neg Q)$. Saranno utili per passare dalle proprietà degli aperti a quelle dei chiusi.

\section{Prodotto cartesiano e relazioni}
% Fonte: Lezione 1, p.~12

\begin{definizione}[Prodotto cartesiano]
Il prodotto cartesiano di $A$ e $B$ è l'insieme delle coppie ordinate con primo elemento in $A$ e secondo elemento in $B$:
\[
  A\times B=\{(a,b):a\in A,\ b\in B\}.
\]
\end{definizione}

\textbf{Coppie ordinate.} Una coppia è \emph{ordinata}: $(a,b)=(a',b')$ se e solo se $a=a'$ e $b=b'$. In particolare, scambiare le componenti modifica in generale la coppia. Per esempio,
\[
  \{1,2\}\times\{a,b\}
   =\{(1,a),(1,b),(2,a),(2,b)\}.
\]

\begin{definizione}[Relazione]
Dati due insiemi non vuoti $A$ e $B$, una relazione tra $A$ e $B$ è un sottoinsieme $R\subseteq A\times B$. Si dice che $a\in A$ e $b\in B$ sono in relazione, e si scrive $a\,R\,b$, quando $(a,b)\in R$. Se $A=B$, si parla di una relazione \emph{su} $A$.
\end{definizione}

\begin{esempio}
Specificare una relazione significa scegliere quali coppie considerare. Su $A=\{1,2,3\}$ la relazione «essere minore o uguale» è il sottoinsieme
\[
  R=\{(1,1),(1,2),(1,3),(2,2),(2,3),(3,3)\}\subseteq A\times A.
\]
\end{esempio}

\section{Relazioni d'ordine e ordine totale}
% Fonte: Lezione 1, p.~13

\begin{definizione}[Relazione d'ordine]
Una relazione $R$ su un insieme non vuoto $A$ si dice \emph{relazione d'ordine} se soddisfa le seguenti proprietà:
\begin{enumerate}[label=(\alph*)]
  \item \textbf{Riflessività:} per ogni $a\in A$, $a\,R\,a$.
  \item \textbf{Antisimmetria:} per ogni $a,b\in A$, $(a\,R\,b\land b\,R\,a)\Rightarrow a=b$.
  \item \textbf{Transitività:} per ogni $a,b,c\in A$, $(a\,R\,b\land b\,R\,c)\Rightarrow a\,R\,c$.
\end{enumerate}
\end{definizione}

\textbf{Antisimmetria e simmetria.} L'antisimmetria non richiede di vietare ogni relazione in entrambi i versi: richiede che ciò avvenga soltanto per elementi uguali. È una condizione diversa dalla simmetria, che richiederebbe invece la validità della relazione in entrambi i versi per ogni coppia in relazione.

\begin{definizione}[Ordine totale]
Un ordine si dice \emph{totale} se ogni coppia di elementi è confrontabile:
\[
  \forall a,b\in A,\qquad a\,R\,b\ \lor\ b\,R\,a.
\]
\end{definizione}

\begin{esempio}
\leavevmode
\begin{itemize}
  \item L'ordine usuale $\le$ su $\N$, $\Z$, $\Q$ o $\R$ è totale.
  \item L'inclusione è un ordine su $\mathcal P(A)$, ma in generale non è totale. Per $A=\{1,2\}$, i sottoinsiemi $\{1\}$ e $\{2\}$ non sono confrontabili: nessuno contiene l'altro. Riflessività, antisimmetria e transitività restano valide.
\end{itemize}
\end{esempio}
